\begin{table*}[!t]
\centering
\normalsize
\setlength{\tabcolsep}{3pt}
\renewcommand{\arraystretch}{1.17}
\begin{tabular}{@{}lccccc@{}}
\toprule
Method & \shortstack{Train\\trajectories} & \shortstack{Train\\pairs} & \shortstack{Optimizer\\time (s)} & \shortstack{NN evals\\/ step} & \shortstack{Inference\\time (s)} \\
\midrule
Increment & 24 & 480 & 280.19 $\pm$ 2.92 & 1.00 & 7.53 $\pm$ 0.83 \\
$\Delta$-Physics & 24 & 480 & 280.19 $\pm$ 2.92 & 1.00 & 8.06 $\pm$ 1.40 \\
Recast-style & 24 & 480 & 280.19 $\pm$ 2.92 & 8.81 & 70.08 $\pm$ 8.50 \\
Cond. F-FNO & 120 & 2400 & 1275.70 $\pm$ 86.78 & 1.00 & 7.38 $\pm$ 0.32 \\
\bottomrule
\end{tabular}
\caption{Data and compute comparison. Inference time is measured on CPU for one viscosity, eight trajectories, and 50 rollout steps at $\nu/\nu_0=2$. One-sided $\Delta$-Physics and Equation-Recast-style reuse the Increment checkpoint and require no additional training. Abbreviations: $\Delta$-Physics denotes One-sided $\Delta$-Physics; Recast-style denotes Equation-Recast-style; Cond. F-FNO denotes Conditioned F-FNO. The first three methods use one training viscosity and 28,370 parameters, with no target-parameter training data; Cond. F-FNO uses five viscosities, target-parameter training data, and 28,386 parameters. Optimizer time excludes data generation.}
\label{tab:data_compute}
\end{table*}
